\medskip

\begin{enumerate}
\item \textbf{Réponse a.}

En effet, la fonction $f$ est dérivable sur $\R_{+}^{*}$, en tant que produit et somme de fonctions dérivables sur cet intervalle. Pour tout $x$ réel strictement positif, on a :

$f'(x)= 1 \times \ln(x) + x \times \dfrac{1}{x} - 1 + 0 = \ln(x) + \dfrac{x}{x} - 1 = \ln(x) +1-1 = \ln (x)$
\item \textbf{Réponse c.}
On a $f(x) = x^2 - x^2\ln (x)$.

$\bullet~~$On a $\displaystyle\lim_{x \to 0} x^2 = 0$ ;

$\bullet~~$On sait (croissances comparées) que  $\displaystyle\lim_{x \to 0} x^2 \ln (x) = 0$ ;

Donc par somme de limites $\displaystyle\lim_{x \to 0}f(x) = 0$.





\item \textbf{Réponse d.}

Quel que soit $x \in \R,\: f(x) = x\left(x^2 - 0,9x - 0,1\right)$.

On a donc $f(x) = 0 \iff \left\{\begin{array}{l c l}
x&=&0\\
x^2 - 0,9x - 0,1&=&0
\end{array}\right.$

$0$ est donc une solution et pour l'équation du second degré le nombre 1 est une racine évidente ; le produit des racines étant égal à $-0,1$ l'autre racine est donc $- 0,1$.

L'équation $f(x) = 0$  a donc trois solutions : $-0,1 \:;\ 0$ et 1.

























			\item \textbf{Réponse c.}

			On dérive la proposition, en pensant bien à utiliser la formule pour dériver une fonction composée ($x\longmapsto H(2x)$ a pour dérivée $x \longmapsto 2H'(2x)$) et on vérifie.

			Si $K_\textbf{c} : x \longmapsto \dfrac{1}{2} H(2x)$, alors,

			$\forall x \in \R, K_\textbf{c}' (x) = \dfrac{1}{2}\times 2H'(2x) = H'(2x) = h(2x)$, car $H$ est une primitive de $h$ sur $\R$

			et donc on a : $\forall x \in \R, K_\textbf{c}' (x) = k(x)$.

			On en déduit donc que $K_\textbf{c}$ est une primitive de $k$ sur $\R$.

			\item \textbf{Réponse b.}

			Dérivons la fonction $f$, qui est dérivable sur $\R$, en tant que produit de fonctions dérivables sur cet ensemble.

			$\forall x \in \R, \quad 1 \times \mathrm{e}^x + x \times \mathrm{e}^x = (1 + x)\mathrm{e}^x$.

			On a donc $f'(1) = (1 + 1)\mathrm{e}^1 = 2\mathrm{e}$.

			Par ailleurs $f(1) = 1\times \mathrm{e}^1 = \mathrm{e}$.

			L'équation réduite de la tangente à la courbe de $f$ au point d'abscisse 1 est donc :
			$\begin{aligned}[t] y &= f'(1) (x - 1) + f(1)\\ &= 2\mathrm{e} (x - 1) + \mathrm{e} \\&= 2\mathrm{e}x - 2\mathrm{e} + \mathrm{e} \\y&= 2\mathrm{e}x - \mathrm{e}\end{aligned}$

			\item \textbf{Réponse d.}

			Résolvons cette inéquation :

			$ \begin{aligned} [t] (0,2)^n < 0,001 &\iff \ln\left(0,2^n\right) < \ln(0,001) \text{\qquad car ln est strictement croissante sur }\R^{*+}\\
			&\iff n\ln(0,02) < \ln(0,001)\\
			&\iff n > \dfrac{\ln(0,001)}{\ln(0,2)} \text{\qquad car } \ln(0,2)< 0
			\end{aligned}$

			On a $\dfrac{\ln(0,001)}{\ln(0,2)} \approx 4,3$, donc les solutions à cette inéquation sont les entiers naturels supérieurs ou égaux à 5.
		\end{enumerate}

\bigskip
